\section{Metodologia}

\subsection{Caracterização da Pesquisa}

Quanto à natureza, a presente pesquisa é classificada como
\textbf{aplicada}, pois busca gerar conhecimento voltado à
solução de um problema real identificado em uma organização
\cite{Gil2002}: o objetivo é analisar o processo atual,
levantar os requisitos, redesenhar e modelar o processo, e
propor a automação de um sistema de gestão de chamados
internos entre áreas, tendo como ponto de entrada um site no
SharePoint e como componentes de automação ferramentas do
ecossistema Microsoft 365.

Quanto à abordagem, a pesquisa é \textbf{qualitativa}, uma
vez que busca compreender as necessidades do contexto
organizacional a partir da perspectiva dos profissionais
diretamente envolvidos no processo, transformando essas
necessidades em requisitos para o redesenho e o
desenvolvimento de uma solução computacional \cite{Gil2002}.

Quanto ao método, a pesquisa caracteriza-se como
\textbf{estudo de caso}, pois analisa um processo específico
--- a gestão de chamados internos entre áreas --- dentro de
uma organização determinada, e desenvolve uma solução
direcionada ao contexto estudado \cite{Gil2002}. Quanto aos
objetivos, a pesquisa é \textbf{exploratória e descritiva}:
exploratória, por investigar a aplicação integrada da
Engenharia de Processos, do BPM e do ecossistema Microsoft
365 na gestão de chamados internos entre áreas; descritiva,
por documentar os requisitos levantados junto aos
participantes, a arquitetura desenvolvida e os resultados
obtidos na avaliação da solução \cite{Santos2023}.

O delineamento adotado é o de \textbf{pesquisa de campo}, uma
vez que a coleta de dados ocorre diretamente no ambiente
organizacional, por meio de um formulário estruturado (Google
Forms) aplicado aos profissionais responsáveis pelo processo
estudado \cite{Gil2002}. Essa escolha permite capturar as
necessidades reais do processo a partir de quem o opera
diariamente, garantindo que os requisitos da solução reflitam
a realidade da organização, além de possibilitar a validação
de todo o fluxo interno da futura automação a partir de
diferentes níveis de atuação dos respondentes.

\subsection{Participantes da Pesquisa}

Os participantes foram selecionados por critério de
intencionalidade, considerando o papel de cada um no processo
de gestão de chamados internos da organização. São três
profissionais da área de Qualidade:

\begin{itemize}
    \item \textbf{Gerente da Área} --- responsável por
    validar a aderência do fluxo proposto aos objetivos
    estratégicos, os benefícios esperados, a viabilidade da
    solução e aprovar o fluxo final;

    \item \textbf{Supervisora de Qualidade} --- responsável
    por validar a aderência do processo, as regras
    operacionais, as necessidades do setor, os impactos na
    rotina da equipe e as prioridades do fluxo;

    \item \textbf{Analista de Qualidade} --- responsável por
    descrever o processo atual, as atividades executadas, as
    dificuldades encontradas, as regras do processo, as
    funcionalidades esperadas da automação e as melhorias
    desejadas.
\end{itemize}

A escolha desses três perfis garante que o levantamento de
requisitos abranja as perspectivas estratégica, operacional e
técnica do processo, permitindo especificar uma solução que
atenda às necessidades de todos os níveis de uso. A divisão
das perguntas do formulário conforme o nível de atuação de
cada participante busca aumentar a qualidade dos requisitos
levantados e validar o fluxo proposto em diferentes níveis
organizacionais.

\subsection{Instrumento de Coleta de Dados}

A coleta de dados foi realizada por meio de um
\textbf{formulário estruturado}, elaborado na ferramenta
\textbf{Google Forms} e aplicado individualmente a cada um dos
três participantes. O formulário foi organizado em seções,
com um bloco de perguntas comuns a todos os respondentes e
blocos específicos direcionados conforme o nível de atuação
do participante (Analista, Supervisora ou Gerente), de modo
que cada perfil respondesse apenas às questões relacionadas ao
seu nível de responsabilidade.

O formulário abordou os seguintes eixos temáticos:
funcionamento atual do processo de comunicação de problemas
entre áreas; principais dificuldades enfrentadas na ausência
de um sistema formalizado; oportunidades de melhoria; regras
de negócio e exceções do processo; atividades executadas
manualmente; campos e informações consideradas essenciais no
registro de um chamado; critérios de priorização e prazos
esperados para resolução; funcionalidades esperadas da
automação; e validação do fluxo proposto em cada nível
organizacional.

As respostas coletadas foram exportadas e analisadas por meio
de análise de conteúdo, com agrupamento em categorias
temáticas que originaram os requisitos funcionais e não
funcionais da solução, posteriormente priorizados pela
técnica MoSCoW e validados pelo orientador.

\subsection{Aspectos Éticos}

A participação no formulário foi voluntária e precedida de
uma seção de apresentação da pesquisa e de consentimento
informado, na qual o participante é esclarecido sobre os
objetivos do levantamento e concorda em responder antes de
prosseguir. As informações coletadas referem-se exclusivamente
ao processo organizacional em análise --- funcionamento do
fluxo, dificuldades, regras de negócio e expectativas quanto à
automação --- não envolvendo dados pessoais sensíveis. Os
participantes são identificados no trabalho apenas pelo
cargo, sem exposição de dados que permitam sua identificação
individual.

Como limitação, destaca-se que os resultados refletem as
características específicas da organização estudada e a
perspectiva dos três participantes selecionados, não podendo
ser generalizados automaticamente para outros contextos
organizacionais, condição inerente ao delineamento de estudo
de caso \cite{Gil2002}.

\subsection{Etapas da Pesquisa}

Em coerência com o ciclo de BPM apresentado na Base Teórica
--- identificação, modelagem, análise, redesenho, automação e
monitoramento --- a pesquisa foi estruturada em seis etapas
sequenciais, apresentadas a seguir.

\subsubsection{Etapa 1 --- Diagnóstico do Processo Atual e
Revisão de Literatura}

Análise do processo existente de comunicação de chamados
internos (processo AS-IS), conduzida em paralelo à revisão de
literatura nas bases Google Scholar, SciELO e Portal de
Periódicos CAPES, com descritores relacionados a engenharia
de processos, BPM, engenharia de requisitos, automação de
processos, plataformas \textit{low-code}, Microsoft 365,
Microsoft List, Power Automate, Power BI e SLA. O objetivo é
construir a fundamentação teórica que embasa o redesenho e o
desenvolvimento da solução proposta.

\subsubsection{Etapa 2 --- Levantamento de Informações via
Google Forms}

Aplicação do formulário estruturado no Google Forms aos três
participantes da área de Qualidade (Analista, Supervisora e
Gerente), com perguntas direcionadas conforme o nível de
atuação de cada perfil. As respostas coletadas foram
analisadas e consolidadas em requisitos funcionais e não
funcionais, priorizados pela técnica MoSCoW (\textit{Must
Have, Should Have, Could Have} e \textit{Won't Have}), e
validados pelo orientador.

\subsubsection{Etapa 3 --- Especificação dos Requisitos}

Transformação das necessidades relatadas pelos participantes
em requisitos funcionais, requisitos não funcionais e regras
de negócio, seguindo a lógica necessidade~$\rightarrow$~análise~
$\rightarrow$~requisito discutida na Base Teórica, com base
nos conceitos de Engenharia de Requisitos.

\subsubsection{Etapa 4 --- Modelagem do Processo Futuro
(TO-BE)}

Redesenho do fluxo de gestão de chamados a partir dos
requisitos especificados, definindo os estados possíveis do
chamado, os critérios de prioridade, o tempo de SLA associado
a cada prioridade e as responsabilidades de cada ator do
processo, segundo a disciplina do BPM.

\subsubsection{Etapa 5 --- Desenvolvimento e Implantação do
Fluxo}

Com base no processo redesenhado, foi desenvolvido o site no
SharePoint responsável por centralizar o acesso à abertura dos
chamados, com os campos, status e critérios de prioridade
definidos na etapa anterior. Em paralelo, foi desenvolvido o
fluxo de automação no Power Automate (geração de protocolo,
notificações e contagem de SLA) e estruturada a Microsoft
List como repositório centralizado. O fluxo foi implantado em
ambiente de testes para verificação do funcionamento
integrado dos componentes, e consolidados os indicadores no
painel gerencial do Power BI (volume de chamados, tempo médio
de resolução, percentual dentro do SLA e distribuição por
área e por prioridade).

\subsubsection{Etapa 6 --- Validação e Consolidação dos
Resultados}

Acompanhamento da validação do processo em uso junto aos
participantes envolvidos, com a consolidação dos dados e
indicadores gerados pelo Power BI. Os resultados são
analisados e confrontados com a hipótese da pesquisa,
discutindo-se à luz da fundamentação teórica a aderência da
proposta às necessidades identificadas nos três níveis
organizacionais consultados.

\begin{table}[H]
\caption{Síntese das etapas da pesquisa}
\label{tab:etapas}
\centering
\begin{tabular}{p{0.5cm} p{3.6cm} p{3.6cm} p{4.4cm}}
\hline
\textbf{N} & \textbf{Etapa} & \textbf{Fonte/Método} &
\textbf{Produto Gerado} \\
\hline
1 & Diagnóstico do processo atual e revisão de literatura &
Análise do processo (AS-IS) e bases científicas &
Fundamentação teórica e caracterização do problema \\
\hline
2 & Levantamento de informações & Formulário Google Forms
com 3 participantes & Necessidades organizacionais
consolidadas \\
\hline
3 & Especificação dos requisitos & Análise de conteúdo e
técnica MoSCoW & Requisitos funcionais e não funcionais \\
\hline
4 & Modelagem do processo futuro (TO-BE) & Redesenho segundo
o ciclo de BPM & Fluxo, estados e regras de negócio
definidos \\
\hline
5 & Desenvolvimento e implantação & SharePoint + Power Apps +
Power Automate + Microsoft List + Power BI & Fluxo implantado
e dashboard lançado \\
\hline
6 & Validação e consolidação dos resultados & Análise
qualitativa junto aos participantes & Resultados e conclusões \\
\hline
\end{tabular}
\raggedright
\par
{\fontsize{10}{12}\selectfont Fonte: Elaborado pelo autor}
\end{table}

\subsection{Cronograma}

O cronograma da pesquisa está estruturado em cinco fases
mensais, distribuídas entre julho e novembro do ano de
execução, conforme descrito na Tabela~\ref{tab:cronograma}.
As Etapas 1 a 4 da pesquisa (diagnóstico, levantamento,
especificação e modelagem) concentram-se no mês de julho; as
Etapas 5 e 6 (desenvolvimento/implantação e
validação/consolidação) correspondem, respectivamente, aos
meses de agosto/setembro e outubro/novembro.

\begin{table}[H]
\caption{Cronograma de execução da pesquisa}
\label{tab:cronograma}
\centering
\begin{tabular}{p{1.5cm} p{3.5cm} p{7cm}}
\hline
\textbf{Mês} & \textbf{Fase} & \textbf{Atividades} \\
\hline
Julho &
Diagnóstico, levantamento, especificação e modelagem &
Análise do processo atual (AS-IS); elaboração e validação do
formulário no Google Forms; aplicação aos três participantes
(Gerente, Supervisora e Analista de Qualidade); análise de
conteúdo das respostas; especificação dos requisitos
funcionais e não funcionais; redesenho do processo futuro
(TO-BE); validação pelo orientador. \\
\hline
Agosto &
Desenvolvimento e implantação do fluxo &
Desenvolvimento do site no SharePoint como ponto de entrada
dos chamados; configuração dos campos obrigatórios, status e
critérios de prioridade; desenvolvimento do fluxo de
automação no Power Automate (geração de protocolo,
notificações e contagem de SLA); estruturação da Microsoft
List; implantação do fluxo em ambiente de testes. \\
\hline
Setembro &
Consolidação, indicadores e lançamento &
Consolidação da integração entre o site no SharePoint e os
demais componentes do ecossistema Microsoft 365; criação dos
indicadores e construção do painel gerencial no Power BI;
lançamento do dashboard do Power BI na primeira semana do
mês; lançamento e início da validação do novo modelo do
processo junto aos participantes. \\
\hline
Outubro &
Validação e consolidação dos resultados &
Acompanhamento da validação do processo em uso; consolidação
dos dados e indicadores gerados; análise preliminar dos
resultados frente à hipótese da pesquisa. \\
\hline
Novembro &
Validação e consolidação dos resultados (continuação) &
Consolidação final dos resultados e da documentação técnica;
redação dos resultados e discussão; revisão final do trabalho
e preparação para a apresentação. \\
\hline
\end{tabular}
\raggedright
\par
{\fontsize{10}{12}\selectfont Fonte: Elaborado pelo autor}
\end{table}
